% !TITLE 楚辞·九歌·少司命
% !AUTHOR 屈原
% !DYNASTY 先秦
% !GENRE 楚辞
% !ORDER 10
% !SOURCE 楚辞·九歌

\begin{poem}
秋兰兮麋芜\textsuperscript{1}，罗生兮堂下\textsuperscript{2}。
绿叶兮素华\textsuperscript{3}，芳菲菲兮袭予\textsuperscript{4}。
夫人自有兮美子\textsuperscript{5}，荪何以兮愁苦\textsuperscript{6}？

秋兰兮青青\textsuperscript{7}，绿叶兮紫茎。
满堂兮美人\textsuperscript{8}，忽独与余兮目成\textsuperscript{9}。
入不言兮出不辞\textsuperscript{10}，乘回风兮载云旗\textsuperscript{11}。
悲莫悲兮生别离\textsuperscript{12}，乐莫乐兮新相知\textsuperscript{13}。

荷衣兮蕙带\textsuperscript{14}，倏而来兮忽而逝\textsuperscript{15}。
夕宿兮帝郊\textsuperscript{16}，君谁须兮云之际\textsuperscript{17}？
与女沐兮咸池\textsuperscript{18}，晞女发兮阳之阿\textsuperscript{19}。
望美人兮未来\textsuperscript{20}，临风怳兮浩歌\textsuperscript{21}。

孔盖兮翠旍\textsuperscript{22}，登九天兮抚彗星\textsuperscript{23}。
竦长剑兮拥幼艾\textsuperscript{24}，荪独宜兮为民正\textsuperscript{25}。
\end{poem}

\begin{annotations}
\notetext{1}{秋兰兮麋芜：秋天的兰草和麋芜（一种香草）。麋芜（mí wú），川芎的幼苗，香草名。开篇以满堂的芬芳香草起兴。}
\notetext{2}{罗生兮堂下：密密地排列生长在堂前。罗，排列、罗列。}
\notetext{3}{素华：白色的花朵。素，白色；华，同“花”。。}
\notetext{4}{芳菲菲兮袭予：芬芳浓郁，扑面而来。袭，侵袭、笼罩。}
\notetext{5}{夫人自有兮美子：人们各自有自己的好儿女。夫（fú），发语词；美子，美好的子女。少司命主宰人间子嗣——人们已经各得子女，神不必忧愁。}
\notetext{6}{荪何以兮愁苦：荪（sūn），一种香草，楚辞中常用作对神或君王的尊称。这里指少司命——你为何要愁苦呢？}
\notetext{7}{青青：通“菁菁”，茂盛的样子。}
\notetext{8}{满堂兮美人：满堂都是美好的人（指来迎神的女巫们）。}
\notetext{9}{忽独与余兮目成：忽然只和我眉目传情。目成，以目传情、两心相许——在满堂美人中，神唯独选中了“我”（主祭的巫者）。}
\notetext{10}{入不言兮出不辞：来的时候不说话，走的时候也不告辞。神的行为神秘莫测，不与人交代。}
\notetext{11}{乘回风兮载云旗：驾着旋风，车上载着云做的旗帜。回风，旋风；云旗，以云为旗。写少司命乘风而来、随风而去的神姿。}
\notetext{12}{悲莫悲兮生别离：人世之悲，莫过于生生的别离。生别离——活着却被分离，比死别更漫长、更折磨。}
\notetext{13}{乐莫乐兮新相知：人世之乐，莫过于新结识一个知己。新相知——初次心灵相通的快乐。}
\notetext{14}{荷衣兮蕙带：荷叶做衣裳，蕙草做腰带。写少司命以香草为衣的仙人之姿。}
\notetext{15}{倏而来兮忽而逝：忽然来了，又忽然去了。倏（shū）、忽，都是疾速的意思。神来去如闪电，不可挽留。}
\notetext{16}{夕宿兮帝郊：傍晚在天帝的郊外住宿。帝郊，天帝之宫的郊野。}
\notetext{17}{君谁须兮云之际：你在云边等待谁？须，等待。“云之际”——少司命站在云的边上，眺望着什么。}
\notetext{18}{与女沐兮咸池：想与你一起在咸池沐浴。女，同“汝”，你；咸池，神话中太阳沐浴的地方。}
\notetext{19}{晞女发兮阳之阿：在日出的山弯里晒干你的头发。晞（xī），晒干；阳之阿（ē），太阳升起的山坳处。}
\notetext{20}{望美人兮未来：遥望美人（少司命），她还没有到来。}
\notetext{21}{临风怳兮浩歌：迎着风，怅然若失，放声高歌。怳（huǎng），恍惚、失落；浩歌，大声歌唱。}
\notetext{22}{孔盖兮翠旍：孔雀羽毛做车盖，翠鸟羽毛做旌旗。孔，孔雀；盖，车盖；翠，翠鸟；旍（jīng），同“旌”，旗。}
\notetext{23}{登九天兮抚彗星：登上九天，轻抚彗星。彗星，即扫帚星。少司命登天抚星的壮丽画面。}
\notetext{24}{竦长剑兮拥幼艾：高高举起长剑，保护着幼小的生灵。竦（sǒng），高举；幼艾，幼儿——少司命是儿童的保护神。}
\notetext{25}{荪独宜兮为民正：只有你（少司命）最适合为民众主持正道。民正，民众的主宰。}
\end{annotations}

\begin{appreciation}
《少司命》和《山鬼》同属《九歌》，都是祭神的乐歌。但《山鬼》写的是等不到的绝望，《少司命》写的却是短暂相会后的怅惘与怀念。如果说山鬼是一个苦苦等待的女子，少司命则更像一场忽然降临、又忽然消失的美梦。

秋兰兮麋芜，罗生兮堂下——满堂香草，芳菲扑面。在一片芬芳中，女神降临了。“满堂兮美人，忽独与余兮目成”——满堂的美人中，她唯独与“我”眉目传情。这一刻的惊喜，写得极为微妙。不是言语，不是动作，只是一个眼神的交会——两个生命在最深的层面完成了理解。但紧接着，“入不言兮出不辞”，神走了。来得没有预告，走得没有告别。那短暂的相知，像彗星划过夜空，还没看清就消失了。

于是情感到此被推上了一个高峰：“悲莫悲兮生别离，乐莫乐兮新相知。”人生最大的悲哀，莫过于活生生的别离；人生最大的快乐，莫过于刚刚结识知己。这两句放在一起，构成了一种残酷的对照：最大的快乐和最大的悲哀之间的距离，就是那一眼目成到神离去之间的距离。

后半部分，是人的追寻。“与女沐兮咸池，晞女发兮阳之阿”，想和你一起在咸池沐浴，在日升之处晒干头发——多么温柔而卑微的愿望。但“望美人兮未来”，她没有来。最后迎风高歌，一个人站在风中失魂落魄。

结尾却是一个壮丽的转身。“竦长剑兮拥幼艾，荪独宜兮为民正”——少司命高举长剑，拥抱着年幼的孩子。她不是谁的专属，她是整个人间的保护神。她必须离开“我”，因为她要去守护所有的孩子。这种从个人情感到神职使命的切换，让全诗的境界豁然开朗：那个瞬间的眼神交汇，已经足够了。神爱众生，一秒的停留已是十足的偏爱。

从这两首九歌的诗篇里，我们看到了一种迥异于《诗经》的篇章所描绘的的世界。同样是写爱情，《诗经》里的爱情是含蓄克制的，而《楚辞》确是直白热烈的。《诗经》所描绘的世界是一个我们熟悉的世界，没有鬼神，充满了烟火气；而《楚辞》描绘的世界却是一个人神共生的世界，祭祀的巫者爱上了女神，山中的精怪在等她的心上人，光怪陆离。但写鬼神，实际上也是在写人。巫者对少司命的爱而不得，与《关雎》、《蒹葭》里对爱情的追求不是一样的么？山鬼苦苦等待她的心上人而不见的悲伤，不就是《氓》里写的“不见复关，泣涕涟涟。既见复关，载笑载言”的心情吗？文学的面貌当然有很多种，但是表达的思想感情却常常相似。这也是我们能跨越千年理解这些诗篇的原因。
\end{appreciation}
